\documentclass[9pt,a4paper]{article}
\usepackage[margin=0.55in]{geometry}
\usepackage[T1]{fontenc}
\usepackage[utf8]{inputenc}
\usepackage{lmodern}
\usepackage{microtype}
\usepackage{graphicx}
\usepackage{booktabs}
\usepackage{xcolor}
\usepackage{hyperref}
\usepackage{caption}
\usepackage{enumitem}
\usepackage{multicol}
\usepackage{titlesec}
\usepackage{float}

\hypersetup{colorlinks=true,linkcolor=blue!50!black,urlcolor=blue!60!black,
  pdftitle={Calculator REST API and CI/CD Pipeline},pdfauthor={Vaishak Prasad Bhat}}
\setlength{\parindent}{0pt}
\setlength{\parskip}{2pt}
\setlist{nosep,leftmargin=*}
\titleformat{\section}{\large\bfseries}{\thesection.}{0.35em}{}
\titleformat{\subsection}{\normalsize\bfseries}{\thesubsection}{0.35em}{}
\titlespacing*{\section}{0pt}{4pt}{2pt}
\titlespacing*{\subsection}{0pt}{3pt}{1pt}
\captionsetup{font=scriptsize,skip=2pt}

\newcommand{\fig}[3]{%
  \begin{minipage}[t]{#2\linewidth}\centering
  \includegraphics[width=\linewidth,height=0.24\textheight,keepaspectratio]{figures/#1}
  \captionof{figure}{#3}
  \end{minipage}}
\newcommand{\widefig}[3]{%
  \begin{figure}[H]\centering
  \includegraphics[width=#2\linewidth,height=0.31\textheight,keepaspectratio]{figures/#1}
  \caption{#3}\end{figure}}

\begin{document}

\begin{center}
{\LARGE\bfseries End-to-End CI/CD Pipeline for a Calculator REST API}\\[-1pt]
{\normalsize CSE 816 -- Software Production Engineering Mini-Project}\\
{\small Vaishak Prasad Bhat \quad|\quad October 2026}\\[2pt]
{\small
\href{https://github.com/vaishak-iiitb/calc_spe}{GitHub repository}
\quad\textbar\quad
\href{https://hub.docker.com/repository/docker/vaishak2005/calculator-api/general}
{Docker Hub image repository}}
\end{center}

\section{Overview}
This project implements a small production-style Calculator REST API with
Java 17, Spring Boot, and Maven, then connects it to GitHub, Jenkins, Docker,
Docker Hub, and Docker Compose. The application performs exact decimal
arithmetic, validates inputs, returns structured HTTP 400 errors, and exposes
an Actuator health probe. The pipeline builds, tests, containerizes, publishes,
and deploys the service automatically. 

\begin{center}
\small
\begin{tabular}{lll}
\toprule
\textbf{Endpoint} & \textbf{Purpose} & \textbf{Result}\\
\midrule
\texttt{/api/v1/add?a=\&b=} & Addition & \(a+b\)\\
\texttt{/api/v1/subtract?a=\&b=} & Subtraction & \(a-b\)\\
\texttt{/api/v1/multiply?a=\&b=} & Multiplication & \(a\times b\)\\
\texttt{/api/v1/divide?a=\&b=} & Division & HTTP 400 for \(b=0\)\\
\texttt{/actuator/health} & Health probe & UP/DOWN\\
\bottomrule
\end{tabular}
\end{center}

\section{Architecture and Deployment}
Docker Compose runs two calculator replicas behind an Nginx reverse proxy.
The proxy publishes port 8081, resolves the service name dynamically, and
routes requests over an internal network. Health checks use
\texttt{/actuator/health}; restart policies and dependency conditions make
the deployment repeatable. Keeping the application containers off the host
port also makes the proxy the single entry point for API requests.

\begin{center}
\fig{screenshot-01.png}{.49}{Compose startup showing two healthy calculator replicas and Nginx.}
\hfill
\fig{screenshot-02.png}{.49}{Addition, division-by-zero, and health responses through Nginx.}
\end{center}

\section{Evidence Summary}
\begin{multicols}{2}
\textbf{Testing:} JUnit 5 unit tests cover every operation, including
division by zero. MockMvc tests exercise real endpoints, validation, and
health. Jenkins publishes JUnit XML results and JaCoCo coverage.

\textbf{Containerization:} A multi-stage Dockerfile builds the JAR in a
Maven/JDK image and runs it in a slim JRE image. The runtime uses a dedicated
non-root user, explicit port exposure, and a health check.

\textbf{Delivery:} Jenkins stages are Checkout, Build, Test, Docker Build,
Push, Deploy, and Post Actions. Docker Hub credentials are injected through
Jenkins Credentials Store rather than committed to source control.

\textbf{Version control:} The repository uses meaningful commits, feature
branches, and a visible non-fast-forward merge of
\texttt{feature/document-git-workflow} into \texttt{main}. The Java/Maven
\texttt{.gitignore} excludes build output and secrets.

\textbf{Reproducibility:} Maven Wrapper is included so the project can be
built without relying on a globally installed Maven version. The Docker and
Compose files keep local deployment commands short and consistent.
\end{multicols}

\newpage
\section{Compose Validation and Load Balancing}
The deployment command \texttt{docker compose up -d --build --wait} builds
the proxy and application images, starts all services, and waits for health.
The screenshots show successful container status and repeated upstream
selection across both calculator replicas. This also checks that Nginx is not
fixed to one container and that both replicas are reachable through the
Compose service name.

\begin{center}
\fig{screenshot-03.png}{.49}{Repeated health checks reaching one calculator upstream.}
\hfill
\fig{screenshot-04.png}{.49}{Repeated health checks reaching the second upstream.}
\end{center}

\widefig{screenshot-05.png}{.82}{Docker run verification: the container runs as the unprivileged user \texttt{app}.}

\section{Quality and Test Results}
The build executes 34 tests: 25 service tests and 9 controller tests, with no
failures or skipped tests. JaCoCo reports 95\% total instruction coverage and
100\% branch coverage for the measured bundle, and the Maven coverage check
enforces the configured threshold.

\begin{center}
\fig{screenshot-06.png}{.49}{JaCoCo package coverage summary.}
\hfill
\fig{screenshot-08.png}{.49}{Jenkins test report: all 34 tests pass.}
\end{center}

\newpage
\section{Jenkins Continuous Delivery}
The declarative \texttt{Jenkinsfile} checks out the source, runs
\texttt{mvn clean package -DskipTests}, runs \texttt{mvn test}, publishes
test and coverage results, builds a build-number-tagged image, pushes it to
Docker Hub, and deploys the Compose stack with health verification.

The pipeline uses Jenkins-managed credentials for Docker Hub authentication.
This keeps login details outside the repository and allows the same pipeline
definition to be reused on another Jenkins installation after credentials and
email settings have been configured. The final deployment stage uses Compose
health checks rather than treating image creation alone as success.

\begin{center}
\fig{screenshot-07.png}{.49}{Build 4: technical stages pass; local SMTP notification is unavailable.}
\hfill
\fig{screenshot-09.png}{.49}{Build 5: all pipeline stages and post actions pass.}
\end{center}

\widefig{screenshot-11.png}{.82}{Email evidence: the successful deployment notification for Jenkins build 17.}

\section{Pipeline Behaviour}
The pipeline is designed to be idempotent: Compose pulls or builds the
requested image, updates services, waits for health, and leaves a running
deployment in a usable state. Notifications report the final build status;
the logged SMTP failure in Build 4 did not conceal the successful technical
stages, while Build 5 and the email in Figure 11 show successful notification
delivery.

\newpage
\section{Published Image and Reproducibility}
The Docker Hub repository contains the published
\texttt{vaishak2005/calculator-api} image with \texttt{latest} and numbered
tags. The public image page is available at
\href{https://hub.docker.com/repository/docker/vaishak2005/calculator-api/general}
{\texttt{hub.docker.com/vaishak2005/calculator-api}}. The image can be run
independently:

% \begin{center}
% \small\texttt{docker run --rm -p 9090:8080 vaishak2005/calculator-api:latest}
% \end{center}

\begin{center}
\fig{screenshot-10.png}{.62}{Docker Hub repository and published image tags.}
\end{center}

\section{Conclusion}
The completed implementation demonstrates the mini-project requirements from
application code through deployment:

\begin{itemize}
\item REST operations, validation, structured errors, and Actuator health;
\item JUnit and MockMvc testing with passing Jenkins results and JaCoCo coverage;
\item secure, multi-stage, non-root Docker packaging;
\item multi-component Compose deployment with Nginx and two replicas;
\item Jenkins automation for build, test, image build, push, deploy, and email;
\item GitHub feature-branch workflow with a meaningful merged history.
\end{itemize}

The final \texttt{main} history contains 19 commits dated from 22--27
October 2026. The application, Docker, Compose, and Jenkins work is organized
as descriptive commits, followed by the feature-branch merge mentioned above.
The source and configuration are available in the
\href{https://github.com/vaishak-iiitb/calc_spe}{GitHub repository}. All
eleven screenshots from \texttt{docs/latex\_report/figures} are included:
Compose status, API responses, upstream distribution, non-root execution,
coverage, Jenkins tests, pipeline outcomes, Docker Hub, and email
notification. Together, they provide a practical record of the checks carried
out while developing and deploying the project.

\end{document}
